% Answers to Chapter 15, translated from sv/back/facit/15.tex.

\facitavsnitt{c15}{Chapter 15}{Complex numbers, part I}
\begin{facit}

\svar{1.1} \sv{a} $u(t) = (3t^2 - 4t + 2)\,\text{V}$ \sv{b} $17\,\text{V}$

\svar{2.1} $U = \sqrt{29}\angle{-0.38}\,\text{rad} \approx 5.39\angle{-0.38}\,\text{rad}$~V,
that is, $\delta_u \approx -21.8^{\circ}$.

\svar{2.2} $I = \frac{5\sqrt{3}}{2} + j2.5 \approx (4.33 + j2.5)\,\text{mA}$

\svar{2.3} $Z \approx 1.08\angle{-0.9}\,\text{rad}$~k$\Omega$
$\approx (0.67 - j0.85)\,\text{k}\Omega$

\svar{3.1} \sv{a} $-1$ \sv{b} $-j$ \sv{c} $1$ \sv{d} $j$ \sv{e} $-j$ \sv{f} $x = \pm j5$

\svar{3.2} \sv{a} Real part and imaginary part: $z_1$: $4$ and $-3$;\enspace $z_2$: $-2$ and
$5$;\enspace $z_3$: $0$ and $6$;\enspace $z_4$: $-7$ and $0$.
\sv{b} $z_3 = j6$ \sv{c} $z_4 = -7$

\svar{3.3} \sv{a} $2 + j7$ \sv{b} $4 - j3$ \sv{c} $3 + j19$ \sv{d} $-7 + j$

\svar{3.4} \sv{a} $-10 + j11$ \sv{b} $13$ \sv{c} $2 + j5$ \sv{d} $j2$

\svar{3.5} \sv{a} $5 - j2$ \sv{b} $-3 + j7$ \sv{c} $25$
\sv{d} $(x + jy)(x - jy) = x^2 - j^2y^2 = x^2 + y^2 = \abs{z}^2$

\svar{3.6} \sv{a} $1 + j3$ \sv{b} $8 - j6$ \sv{c} $1.2 - j1.6$
\sv{d} The product of the denominator and its conjugate is real, and with a real denominator the
real and imaginary parts can be divided separately.

\svar{3.7} \sv{a} $5$ \sv{b} $13$ \sv{c} $7$ \sv{d} $8$ \sv{e} $\sqrt{2} \approx 1.41$

\svar{3.8} \sv{a} $45^{\circ}$, quadrant I \sv{b} $135^{\circ}$, quadrant II
\sv{c} $225^{\circ}$ (or $-135^{\circ}$), quadrant III \sv{d} $-45^{\circ}$, quadrant IV
\sv{e} $53.1^{\circ}$, quadrant I

\svar{3.9} \sv{a} $10\angle 53.1^{\circ}$ \sv{b} $5\angle 126.9^{\circ}$
\sv{c} $\sqrt{2}\angle 225^{\circ} \approx 1.41\angle 225^{\circ}$ \sv{d} $5\angle{-90^{\circ}}$

\svar{3.10} \sv{a} $8.66 + j5$ \sv{b} $j4$ \sv{c} $-6$ \sv{d} $1 - j1.73$
\sv{e} $5.66 + j5.66$

\svar{3.11} \sv{a} $13\,\text{V}$ \sv{b} $22.6^{\circ}$ \sv{c} $13\angle 22.6^{\circ}\,\text{V}$
\sv{d} How far the voltage is shifted in time relative to the reference, for example the current.

\svar{3.12} \sv{a} $Z = (30 + j40)\,\Omega$ \sv{b} $50\,\Omega$ \sv{c} $53.1^{\circ}$
\sv{d} Inductive: the imaginary part is positive.

\svar{3.13} \sv{a} $\abs{U} = 10\,\text{V}$, $\abs{Z} = 5\,\Omega$
\sv{b} $I = (1.2 - j1.6)\,\text{A}$ \sv{c} $\abs{I} = \frac{10}{5} = 2\,\text{A}$
\sv{d} $-53.1^{\circ}$: the current lags the voltage by $53.1^{\circ}$, as in an inductive
circuit.

\svar{3.14} \sv{a} $x = \pm j3$ \sv{b} $x = -2 \pm j3$ \sv{c} $x = 1 \pm j2$
\sv{d} They are each other's conjugates: the same real part and opposite imaginary parts.

\svar{3.15} \sv{a} $8\,\Omega$. The imaginary parts cancel, so the series connection is purely
resistive. \sv{b} Both are $5\,\Omega$. \sv{c} $25\,\Omega^2$ \sv{d} $3.125\,\Omega$

\svar{3.16} \sv{a} False: $j^2 = -1$.
\sv{b} False: $\abs{z}$ is a length, always $\geq 0$.
\sv{c} True: the number lies on the imaginary axis.
\sv{d} True: the product is $a^2 + b^2$.
\sv{e} False: $\arctan$ only gives angles between $-90^{\circ}$ and $90^{\circ}$; in quadrants II
and III, $180^{\circ}$ is added.
\sv{f} False: addition is easiest in rectangular form.

\svar[Test yourself]{T} \sv{1} $\abs{z} = 5$ and $\delta \approx 0.927\,\text{rad}$, so
$z = 5\angle 0.927$. \sv{2} $5 + j$

\end{facit}
